\hojaejercicio{Lista 2}





%ejercicioN(1)
\ejercicio{Considera una homografía de $\mathbb{P}_{\mathbb{C}}^{2}$ que transforma los puntos $(1:0:-1)$ y $(1:1:-1)$ en los puntos $(0:1:-1)$ y $(0:2:1)$ y en alguna carta afín es una traslación.\\
a) Halla la recta de infinito de la carta afín anterior.\\
b) Obtén las ecuaciones de la homografía.}{}

%ejercicioN(2)
\ejercicio{Obtén las ecuaciones de una homografía de $\mathbb{P}_{\mathbb{R}}^{2}$ que:
\begin{itemize}
    \item Deja fijos los puntos $(1:1:0), (0:0:1)$ y $(-1:1:0)$.
    \item Transforma el punto $(-1:1:-1)$ en $(-1:1:1)$.
    \item En alguna carta afín es una homotecia.
\end{itemize}
¿Hay más de una homografía con las anteriores características?}{
    \section{Definición de la Referencia de Puntos Fijos}
Consideramos los tres puntos fijos dados como la base de un nuevo sistema de referencia $\mathcal{B} = \{P_1, P_2, P_3\}$:
$$ P_1 = (1:1:0), \quad P_2 = (0:0:1), \quad P_3 = (-1:1:0) $$

La matriz de cambio de base $P$ (que transforma coordenadas de $\mathcal{B}$ a la base canónica $\mathcal{C}$) se construye colocando los vectores en columnas:
$$ P = \begin{pmatrix} 1 & 0 & -1 \\ 1 & 0 & 1 \\ 0 & 1 & 0 \end{pmatrix} $$

La matriz inversa (necesaria para volver a la canónica) es:
$$ P^{-1} = \begin{pmatrix} 0.5 & 0.5 & 0 \\ 0 & 0 & 1 \\ -0.5 & 0.5 & 0 \end{pmatrix} $$


Expresamos el punto original $Q(-1:1:-1)$ y su imagen $Q'(-1:1:1)$ en la base $\mathcal{B}$ resolviendo el sistema $X_{\mathcal{B}} = P^{-1} \cdot X_{\mathcal{C}}$:

\begin{itemize}
    \item \textbf{Punto $Q$:} $Q_{\mathcal{B}} = P^{-1} \begin{pmatrix} -1 \\ 1 \\ -1 \end{pmatrix} = \begin{pmatrix} 0 \\ -1 \\ 1 \end{pmatrix}$
    \item \textbf{Punto $Q'$:} $Q'_{\mathcal{B}} = P^{-1} \begin{pmatrix} -1 \\ 1 \\ 1 \end{pmatrix} = \begin{pmatrix} 0 \\ 1 \\ 1 \end{pmatrix}$
\end{itemize}


Buscamos una homología donde el Centro tenga autovalor $1$ y la recta de puntos fijos tenga autovalor $k$ (con $k \neq 1$).

Caso 1: Centro $P_1$ (Matriz $M_{\mathcal{B}} = \text{diag}(1, k, k)$)
$$ \begin{pmatrix} 1 & 0 & 0 \\ 0 & k & 0 \\ 0 & 0 & k \end{pmatrix} \begin{pmatrix} 0 \\ -1 \\ 1 \end{pmatrix} = \begin{pmatrix} 0 \\ -k \\ k \end{pmatrix} \sim (0:-1:1) \neq (0:1:1) $$
No existe solución para $k$. \textbf{Caso Imposible.}

Caso 2: Centro $P_2$ (Matriz $M_{\mathcal{B}} = \text{diag}(k, 1, k)$)
$$ \begin{pmatrix} k & 0 & 0 \\ 0 & 1 & 0 \\ 0 & 0 & k \end{pmatrix} \begin{pmatrix} 0 \\ -1 \\ 1 \end{pmatrix} = \begin{pmatrix} 0 \\ -1 \\ k \end{pmatrix} \sim \begin{pmatrix} 0 \\ 1 \\ 1 \end{pmatrix} \implies \frac{-1}{1} = \frac{k}{1} \implies \mathbf{k = -1} $$

Caso 3: Centro $P_3$ (Matriz $M_{\mathcal{B}} = \text{diag}(k, k, 1)$)
$$ \begin{pmatrix} k & 0 & 0 \\ 0 & k & 0 \\ 0 & 0 & 1 \end{pmatrix} \begin{pmatrix} 0 \\ -1 \\ 1 \end{pmatrix} = \begin{pmatrix} 0 \\ -k \\ 1 \end{pmatrix} \sim \begin{pmatrix} 0 \\ 1 \\ 1 \end{pmatrix} \implies \frac{-k}{1} = \frac{1}{1} \implies \mathbf{k = -1} $$


Para obtener las ecuaciones finales, calculamos $M_{\mathcal{C}} = P \cdot M_{\mathcal{B}} \cdot P^{-1}$.

Solución 1 (Centro $P_2$)
Utilizando $M_{\mathcal{B}} = \text{diag}(-1, 1, -1)$:
$$ M_{\mathcal{C}} = \begin{pmatrix} 1 & 0 & -1 \\ 1 & 0 & 1 \\ 0 & 1 & 0 \end{pmatrix} \begin{pmatrix} -1 & 0 & 0 \\ 0 & 1 & 0 \\ 0 & 0 & -1 \end{pmatrix} \begin{pmatrix} 0.5 & 0.5 & 0 \\ 0 & 0 & 1 \\ -0.5 & 0.5 & 0 \end{pmatrix} = \begin{pmatrix} -1 & 0 & 0 \\ 0 & -1 & 0 \\ 0 & 0 & 1 \end{pmatrix} \sim \begin{pmatrix} 1 & 0 & 0 \\ 0 & 1 & 0 \\ 0 & 0 & -1 \end{pmatrix} $$
\textbf{Ecuaciones:} $x_0' = x_0, \quad x_1' = x_1, \quad x_2' = -x_2$.

Solución 2 (Centro $P_3$)
Utilizando $M_{\mathcal{B}} = \text{diag}(-1, -1, 1)$:
$$ M_{\mathcal{C}} = \begin{pmatrix} 1 & 0 & -1 \\ 1 & 0 & 1 \\ 0 & 1 & 0 \end{pmatrix} \begin{pmatrix} -1 & 0 & 0 \\ 0 & -1 & 0 \\ 0 & 0 & 1 \end{pmatrix} \begin{pmatrix} 0.5 & 0.5 & 0 \\ 0 & 0 & 1 \\ -0.5 & 0.5 & 0 \end{pmatrix} = \begin{pmatrix} 0 & -1 & 0 \\ -1 & 0 & 0 \\ 0 & 0 & -1 \end{pmatrix} \sim \begin{pmatrix} 0 & 1 & 0 \\ 1 & 0 & 0 \\ 0 & 0 & 1 \end{pmatrix} $$
\textbf{Ecuaciones:} $x_0' = x_1, \quad x_1' = x_0, \quad x_2' = x_2$.
}

%ejercicioN(3)
\ejercicio{Sea la proyección cónica de $\mathbb{P}_{\mathbb{R}}^{3}$ sobre el plano $x_{0}=0$ de centro el punto $(1:0:1:0)$. Halla unas ecuaciones de la proyección paralela (afín) inducida por en el complementario del plano $x_{0}-x_{2}=0$.}{
    Para simplificar el problema, definimos una nueva base $B = \{ \vec{v}_0, \vec{v}_1, \vec{v}_2, \vec{v}_3 \}$ de $\mathbb{R}^4$ tal que el centro de proyección $C$ y el plano del infinito $H_\infty$ queden alineados con los ejes:
\begin{itemize}
    \item $\vec{v}_0 = (1, 0, 1, 0)$ (Centro de proyección $C$).
    \item $\vec{v}_1 = (0, 1, 0, 0) = \vec{e}_2$.
    \item $\vec{v}_2 = (0, 0, 1, 0) = \vec{e}_3$.
    \item $\vec{v}_3 = (0, 0, 0, 1) = \vec{e}_4$.
\end{itemize}

La relación entre las coordenadas originales $(x_i)$ y las nuevas coordenadas $(X_i)_{\mathcal{R}}$ viene dada por:
\[
\begin{pmatrix} x_0 \\ x_1 \\ x_2 \\ x_3 \end{pmatrix} = 
\begin{pmatrix} 
1 & 0 & 0 & 0 \\ 
0 & 1 & 0 & 0 \\ 
1 & 0 & 1 & 0 \\ 
0 & 0 & 0 & 1 
\end{pmatrix} 
\begin{pmatrix} X_0 \\ X_1 \\ X_2 \\ X_3 \end{pmatrix} 
\implies 
\begin{cases} 
x_0 = X_0 \\ 
x_1 = X_1 \\ 
x_2 = X_0 + X_2 \\ 
x_3 = X_3 
\end{cases}
\]


\begin{itemize}
    \item \textbf{Plano del infinito:} El enunciado define el espacio afín como el complementario de $x_0 - x_2 = 0$. Sustituyendo las expresiones obtenidas:
    \[ X_0 - (X_0 + X_2) = 0 \implies \mathbf{X_2 = 0} \]
    En esta referencia, la carta afín se obtiene tomando $X_2 = 1$.
    
    \item \textbf{Centro de proyección:} $C = (1 : 0 : 0 : 0)_{\mathcal{R}}$. Al pertenecer al plano $X_2=0$, el centro es un punto del infinito, lo que garantiza que la proyección inducida es \textbf{paralela}.
    
    \item \textbf{Plano imagen:} $x_0 = 0 \implies \mathbf{X_0 = 0}$.
\end{itemize}


La proyección cónica $\pi$ desde el centro $(1:0:0:0)$ sobre el plano $X_0=0$ en coordenadas homogéneas es:
\[ \pi(X_0 : X_1 : X_2 : X_3) = (0 : X_1 : X_2 : X_3) \]

Pasamos al espacio afín tomando la carta $\mathbf{X_2 = 1}$. Un punto genérico de $\mathbb{A}^3$ tiene coordenadas:
\[ P = (y_1 : y_2 : 1 : y_3)_{\mathcal{R}} \quad \text{donde } y_1 = \frac{X_0}{X_2}, y_2 = \frac{X_1}{X_2}, y_3 = \frac{X_3}{X_2} \]

Aplicamos la proyección evaluando para $X_2=1$:
\[ \pi(P) = (0 : y_2 : 1 : y_3)_{\mathcal{R}} \]

Las coordenadas de la imagen $(z_1, z_2)$ en el plano afín destino (subespacio de $\mathbb{A}^3$ donde $X_0=0$) son:
\[ \begin{cases} z_1 = y_2 \\ z_2 = y_3 \end{cases} \]

\textbf{Conclusión:} La proyección afín inducida es la proyección paralela que mantiene las coordenadas $y_2, y_3$ y anula la coordenada $y_1$ (dirección del centro de proyección).}

%ejercicioN(4)
\ejercicio{Sea $f$ la homografía de $\mathbb{P}_{\mathbb{R}}^{3}$ de ecuaciones:
\[ \begin{cases} \lambda x_{0}^{\prime}=x_{0}+x_{1} \\ \lambda x_{1}^{\prime}=2x_{1}+x_{2} \\ \lambda x_{2}^{\prime}=2x_{2}+x_{3} \\ \lambda x_{3}^{\prime}=x_{3} \end{cases} \]
a) Halla sus puntos fijos y sus planos invariantes.\\
b) Halla una referencia proyectiva $R$ del plano $x_{2}+x_{3}=0$ y obtén la matriz de la restricción de $f$ a dicho plano en la referencia $R$.\\
c) Calcula las rectas invariantes contenidas en los planos $x_{2}+x_{3}=0$ y $x_{3}=0$.\\
Unos apartados más exigentes:\\
d) Prueba que las rectas obtenidas en el apartado anterior son todas las rectas invariantes por $f$. (Pista: si una recta corta en un punto a un plano invariante lo hace en un punto fijo, y si dos rectas coplanarias son invariantes entonces el plano que generan también lo es.)\\
e) Halla unas ecuaciones de la aplicación afín inducida en el plano afín $\{x_{3}=0\}\setminus\{x_{2}=0\}$.}{
    A)
    Los Puntos Fijos son:
Los puntos fijos de la homografía corresponden a las direcciones de los autovectores de la matriz $A$. Calculamos el polinomio característico:
\[ p(\lambda) = \det(A - \lambda I) = (1-\lambda)^2(2-\lambda)^2 \]
Los autovalores son $\lambda_1 = 1$ (doble) y $\lambda_2 = 2$ (doble).

\begin{itemize}
    \item \textbf{Para $\lambda = 1$:} Resolvemos $(A-I)\vec{v} = \vec{0}$:
    \[ \begin{pmatrix} 0 & 1 & 0 & 0 \\ 0 & 1 & 1 & 0 \\ 0 & 0 & 1 & 1 \\ 0 & 0 & 0 & 0 \end{pmatrix} \begin{pmatrix} x_0 \\ x_1 \\ x_2 \\ x_3 \end{pmatrix} = \begin{pmatrix} 0 \\ 0 \\ 0 \\ 0 \end{pmatrix} \implies x_1=0, x_2=0, x_3=0 \]
    El subespacio propio es $V_1 = \langle (1,0,0,0) \rangle$. El punto fijo es \textbf{$P_1(1:0:0:0)$}.

    \item \textbf{Para $\lambda = 2$:} Resolvemos $(A-2I)\vec{v} = \vec{0}$:
    \[ \begin{pmatrix} -1 & 1 & 0 & 0 \\ 0 & 0 & 1 & 0 \\ 0 & 0 & 0 & 1 \\ 0 & 0 & 0 & -1 \end{pmatrix} \begin{pmatrix} x_0 \\ x_1 \\ x_2 \\ x_3 \end{pmatrix} = \begin{pmatrix} 0 \\ 0 \\ 0 \\ 0 \end{pmatrix} \implies x_0=x_1, x_2=0, x_3=0 \]
    El subespacio propio es $V_2 = \langle (1,1,0,0) \rangle$. El punto fijo es \textbf{$P_2(1:1:0:0)$}.
\end{itemize}

Los Planos Invariantes son:
Un plano $u_0x_0 + u_1x_1 + u_2x_2 + u_3x_3 = 0$ es invariante si el vector de sus coeficientes $u = (u_0, u_1, u_2, u_3)^T$ es un autovector de la matriz traspuesta $A^T$:
\[ A^T = \begin{pmatrix} 1 & 0 & 0 & 0 \\ 1 & 2 & 0 & 0 \\ 0 & 1 & 2 & 0 \\ 0 & 0 & 1 & 1 \end{pmatrix} \]

\begin{itemize}
    \item \textbf{Para $\lambda = 1$:} $(A^T-I)u = 0 \implies u_0+u_1=0, u_1+u_2=0, u_2=0$. La solución es $u = (0,0,0,1)$. 
    El plano invariante es \textbf{$\pi_1: x_3 = 0$}.
    \item \textbf{Para $\lambda = 2$:} $(A^T-2I)u = 0 \implies u_0=0, u_1=0, u_2-u_3=0$. La solución es $u = (0,0,1,1)$. 
    El plano invariante es \textbf{$\pi_2: x_2 + x_3 = 0$}.
\end{itemize}

B)
Definimos una referencia proyectiva $\mathcal{R}$ en el plano $x_2+x_3=0$ usando los vectores:
\[ \vec{v}_1 = (1,0,0,0), \quad \vec{v}_2 = (0,1,0,0), \quad \vec{v}_3 = (0,0,1,-1) \]
Calculamos la imagen de estos vectores por $A$ y las expresamos en términos de la misma base $\mathcal{R}$ para obtener la matriz restringida $M_{\mathcal{R}}$:
\begin{enumerate}
    \item $A(1,0,0,0) = (1,0,0,0) = \mathbf{1}\vec{v}_1 + \mathbf{0}\vec{v}_2 + \mathbf{0}\vec{v}_3$
    \item $A(0,1,0,0) = (1,2,0,0) = \mathbf{1}\vec{v}_1 + \mathbf{2}\vec{v}_2 + \mathbf{0}\vec{v}_3$
    \item $A(0,0,1,-1) = A(0,0,1,0) - A(0,0,0,1) = (0,1,2,0) - (0,0,1,1) = (0,1,1,-1) = \mathbf{0}\vec{v}_1 + \mathbf{1}\vec{v}_2 + \mathbf{1}\vec{v}_3$
\end{enumerate}
La matriz de la restricción es:
\[ M_{\mathcal{R}} = \begin{pmatrix} 1 & 1 & 0 \\ 0 & 2 & 1 \\ 0 & 0 & 1 \end{pmatrix} \]

C)
Una recta contenida en un plano invariante es invariante si su vector dual en la restricción es autovector de $M_{restr}^T$.

En el plano $x_2+x_3=0$ (Matriz $M_{\mathcal{R}}$)
Analizamos los autovectores de $M_{\mathcal{R}}^T = \begin{pmatrix} 1 & 0 & 0 \\ 1 & 2 & 0 \\ 0 & 1 & 1 \end{pmatrix}$:
\begin{itemize}
    \item $\lambda=1$: Autovector $(0,0,1)_{\mathcal{R}}$. Representa la recta $z_{\mathcal{R}}=0$. En el espacio original, esta es la recta que une $\vec{v}_1$ y $\vec{v}_2$: \textbf{$x_2=0, x_3=0$}.
    \item $\lambda=2$: Autovector $(0,1,1)_{\mathcal{R}}$. Representa la recta $y_{\mathcal{R}}+z_{\mathcal{R}}=0$. En el espacio original: \textbf{$x_1+x_2=0, x_2+x_3=0$}.
\end{itemize}

En el plano $x_3=0$:
La matriz de la restricción en la base $\{(1,0,0,0), (0,1,0,0), (0,0,1,0)\}$ es:
\[ N = \begin{pmatrix} 1 & 1 & 0 \\ 0 & 2 & 1 \\ 0 & 0 & 2 \end{pmatrix} \implies N^T = \begin{pmatrix} 1 & 0 & 0 \\ 1 & 2 & 0 \\ 0 & 1 & 2 \end{pmatrix} \]
\begin{itemize}
    \item $\lambda=1$: Autovector $(1,-1,1)$. Recta: \textbf{$x_0 - x_1 + x_2 = 0$ (con $x_3=0$)}.
    \item $\lambda=2$: Autovector $(0,0,1)$. Recta: \textbf{$x_2 = 0$ (con $x_3=0$)}.
\end{itemize}
D)
\begin{enumerate}
    \item \textbf{Anclaje a los puntos fijos (Pista 1):} 
    Sea $r$ una recta invariante por la homografía $f$. En el espacio proyectivo $\mathbb{P}^3$, toda recta tiene intersección no vacía con cualquier plano. Consideremos los planos invariantes $\pi_1$ y $\pi_2$ hallados en el apartado (a).
    
    Por la \textit{Pista 1}, los puntos de intersección $Q_1 = r \cap \pi_1$ y $Q_2 = r \cap \pi_2$ deben ser necesariamente puntos fijos de la homografía. Dado que los únicos puntos fijos son $P_1(1:0:0:0)$ y $P_2(1:1:0:0)$, existen dos posibilidades para $r$:
    \begin{itemize}
        \item Que $r$ pase por dos puntos fijos distintos ($P_1$ y $P_2$). Esta es la recta $r_1$, ya calculada.
        \item Que $r$ pase por un único punto fijo (por ejemplo, $Q_1 = Q_2 = P_1$). En este caso, la recta interseca a ambos planos en su punto de corte común.
    \end{itemize}

    \item \textbf{Restricción a los planos existentes (Pista 2):}
    Supongamos que existiera una recta invariante $r$ que pasara por un punto fijo (pongamos $P_1$) pero que no estuviera contenida ni en $\pi_1$ ni en $\pi_2$.
    
    Tomemos una recta invariante conocida que también pase por $P_1$, como es $r_1$. Al cortarse ambas en $P_1$, las rectas $r$ y $r_1$ son coplanarias. Según la \textit{Pista 2}, el plano $\alpha$ generado por $r$ y $r_1$ debe ser también un plano invariante.
    
    \item \textbf{Contradicción con el estudio de planos:}
    En el apartado (a) se determinó, mediante el cálculo de los autovectores de la matriz traspuesta $A^T$, que los únicos planos invariantes de la transformación son $\pi_1$ y $\pi_2$. 
    
    Si existiera la recta $r$ fuera de estos planos, el plano $\alpha$ sería un tercer plano invariante distinto de los anteriores, o bien formaría parte de un haz de planos invariantes. Ambas situaciones contradicen el hecho de que el subespacio propio de $A^T$ para cada autovalor tiene dimensión 1, lo que limita los planos invariantes a solo dos.
\end{enumerate}
E)
Consideramos la restricción de la homografía al plano invariante $x_3 = 0$. La matriz reducida que actúa sobre las coordenadas $(x_0 : x_1 : x_2)$ es:
\[
M = \begin{pmatrix} 1 & 1 & 0 \\ 0 & 2 & 1 \\ 0 & 0 & 2 \end{pmatrix}
\]

Para obtener la aplicación afín, observamos el plano desde la carta afín donde la recta del infinito es $x_2 = 0$. Esto implica fijar $x_2 = 1$ y definir nuestras coordenadas afines como $y_1 = x_0$ e $y_2 = x_1$.

\begin{enumerate}
    \item \textbf{Transformación en coordenadas homogéneas:}
    Tomamos un punto genérico del plano afín expresado como $(y_1, y_2, 1)$ y aplicamos la matriz reducida:
    \[
    \lambda \begin{pmatrix} x'_0 \\ x'_1 \\ x'_2 \end{pmatrix} = 
    \begin{pmatrix} 1 & 1 & 0 \\ 0 & 2 & 1 \\ 0 & 0 & 2 \end{pmatrix}
    \begin{pmatrix} y_1 \\ y_2 \\ 1 \end{pmatrix} = 
    \begin{pmatrix} y_1 + y_2 \\ 2y_2 + 1 \\ 2 \end{pmatrix}
    \]

    \item \textbf{Normalización por proporcionalidad:}
    Como el resultado proyectivo es $(y_1 + y_2 : 2y_2 + 1 : 2)$, para obtener las coordenadas en nuestro plano "afín" (donde la tercera componente debe ser $1$), debemos dividir todo el vector por el valor de la coordenada de referencia, que es $2$. 
    
    Este paso de normalización asegura que el punto imagen se mantenga en la misma carta afín $x_2=1$:
    \[
    y'_1 = \frac{y_1 + y_2}{2}, \quad y'_2 = \frac{2y_2 + 1}{2}
    \]

    \item \textbf{Ecuaciones de la aplicación afín:}
    Separando los términos, obtenemos la expresión final de la aplicación:
    \[
    \begin{cases}
    y'_1 = \frac{1}{2}y_1 + \frac{1}{2}y_2 \\
    y'_2 = y_2 + \frac{1}{2}
    \end{cases}
    \]
\end{enumerate}
}

%ejercicioN(5)
\ejercicio{Sea $g$ la homografía de $\mathbb{P}_{\mathbb{R}}^{3}$ de ecuaciones:
\[ \begin{cases} \lambda x_{0}^{\prime}=-x_{0}+2x_{1} \\ \lambda x_{1}^{\prime}=x_{1} \\ \lambda x_{2}^{\prime}=x_{1}-x_{2} \\ \lambda x_{3}^{\prime}=-2x_{1}+3x_{2}+2x_{3} \end{cases} \]
a) Halla sus puntos fijos y sus planos invariantes.\\
b) Comprueba dentro del conjunto de planos invariantes se encuentra un haz de planos cuya base es la recta generada por $(2:2:1:1)$ y $(0:0:0:1)$.\\
Unos apartados más exigentes:\\
c) Halla todas las rectas invariantes por $g$.\\
d) Determina la recta $r$ del plano $\{x_{1}=0\}$ tal que $g$ induce una homotecia en el plano afín $\{x_{1}=0\}\setminus r$.\\
e) Halla el centro y la razón de la homotecia.}{
    Dada la homografía $g$ de $\mathbb{P}^3_{\mathbb{R}}$, extraemos su matriz asociada $A$:
\[
A = \begin{pmatrix} 
-1 & 2 & 0 & 0 \\ 
0 & 1 & 0 & 0 \\ 
0 & 1 & -1 & 0 \\ 
0 & -2 & 3 & 2 
\end{pmatrix}
\]

a) 
1. Cálculo de Autovalores:
Resolvemos el polinomio característico $|A - \lambda I| = 0$:
\[
(-1-\lambda)^2(1-\lambda)(2-\lambda) = 0 \implies \lambda_1 = -1 \text{ (doble)}, \lambda_2 = 1, \lambda_3 = 2
\]

2. Puntos Fijos (Autovectores de $A$):
\begin{itemize}
    \item \textbf{Para $\lambda = -1$}: El sistema $(A+I)\vec{x} = \vec{0}$ da $x_1 = 0$ y $x_2 = -x_3$. La variable $x_0$ es libre. Obtenemos la \textbf{recta de puntos fijos} generada por $P_1(1:0:0:0)$ y $P_2(0:0:1:-1)$.
    \item \textbf{Para $\lambda = 1$}: El sistema $(A-I)\vec{x} = \vec{0}$ da el punto \textbf{$P_3(2:2:1:1)$}.
    \item \textbf{Para $\lambda = 2$}: El sistema $(A-2I)\vec{x} = \vec{0}$ da el punto \textbf{$P_4(0:0:0:1)$}.
\end{itemize}
3. Planos Invariantes (Autovectores de $A^T$):
Buscamos los planos $\pi \equiv u_0x_0 + u_1x_1 + u_2x_2 + u_3x_3 = 0$ tales que $(A^T - \lambda I)\vec{u} = \vec{0}$.
\begin{itemize}
    \item \textbf{Para $\lambda = -1$}: El sistema $(A^T+I)\vec{u} = \vec{0}$ impone las condiciones:
    \[ u_3 = 0 \quad \text{y} \quad 2u_0 + 2u_1 + u_2 = 0 \implies u_2 = -2u_0 - 2u_1 \]
    Esto define un \textbf{haz de planos} cuya ecuación general es:
    \[ \pi_{u_0, u_1} \equiv u_0x_0 + u_1x_1 + (-2u_0 - 2u_1)x_2 = 0 \]
    \item \textbf{Para $\lambda = 1$}: Obtenemos el plano \textbf{$\pi_3 \equiv x_1 = 0$}.
    \item \textbf{Para $\lambda = 2$}: Obtenemos el plano \textbf{$\pi_4 \equiv -x_1 + x_2 + x_3 = 0$}.
\end{itemize}

b) 

El enunciado nos pide comprobar si dentro del conjunto de planos invariantes se encuentra un haz cuya base es la recta generada por $Q_1(2:2:1:1)$ y $Q_2(0:0:0:1)$. 

Para demostrarlo, comprobaremos que ambos puntos pertenecen a \textbf{todos} los planos del haz asociado a $\lambda = -1$:
\[ \pi \equiv u_0x_0 + u_1x_1 - (2u_0 + 2u_1)x_2 = 0 \]

\begin{enumerate}
    \item Comprobación para $Q_2(0:0:0:1)$:
    Sustituimos las coordenadas $(x_0, x_1, x_2, x_3) = (0, 0, 0, 1)$ en la ecuación del haz:
    \[ u_0(0) + u_1(0) - (2u_0 + 2u_1)(0) = 0 \]
    Se cumple $0=0$ para cualquier valor de $u_0$ y $u_1$. Por tanto, $Q_2$ está en todos los planos del haz.
    
    \item Comprobación para $Q_1(2:2:1:1)$:
    Sustituimos las coordenadas $(x_0, x_1, x_2, x_3) = (2, 2, 1, 1)$ en la ecuación:
    \[ u_0(2) + u_1(2) - (2u_0 + 2u_1)(1) = 2u_0 + 2u_1 - 2u_0 - 2u_1 = 0 \]
    Se cumple $0=0$ independientemente de los coeficientes $u_0, u_1$. Por tanto, $Q_1$ también pertenece a todos los planos del haz.
\end{enumerate}

\textbf{Conclusión:} Al estar $Q_1$ y $Q_2$ contenidos en todos los planos invariantes asociados al autovalor $\lambda = -1$, la recta que los une es necesariamente la base (eje) de dicho haz.
c)
Las rectas invariantes de $g$ son:

\begin{enumerate}
    \item \textbf{La recta de puntos fijos:} 
    $r_{12} = \langle (1:0:0:0), (0:0:1:-1) \rangle$. Es invariante punto a punto pues está asociada al autovalor $\lambda = -1$.
    
    \item \textbf{La recta que une los puntos fijos con autovalores $\lambda=1$ y $\lambda=2$:}
    $r_{34} = \langle (2:2:1:1), (0:0:0:1) \rangle$. Esta recta es invariante.
    
    \item \textbf{El conjunto de rectas que unen $P_3$ con cualquier punto de $L_{12}$:}
    Estas rectas forman un haz de rectas con centro $P_3$ contenido en el plano invariante $\pi_4 \equiv -x_1 + x_2 + x_3 = 0$.
    
    \item \textbf{El conjunto de rectas que unen $P_4$ con cualquier punto de $L_{12}$:}
    Estas rectas forman un haz de rectas con centro $P_4$ contenido en el plano invariante $\pi_3 \equiv x_1 = 0$.
    \\
    Usando un resultado visto en clase, sabemos que en un cuerpo algebraicamente cerrado, los subespacios invariantes contienen un punto fijo y están contenidos por un hiperplano invariante. Así, como en nuestro caso, todos autovalores son reales y los autovectores también. Por lo tanto, los puntos fijos de la aplicación en $\mathbb{C}$   son de $P(\mathbb{R}^4)$  y los hiperplanos invariantes tambien. Así, podemos asegurar que cubrimos todas las posibilidades.
\end{enumerate}
d)

Para que la homografía $g$ induzca una homotecia en el plano afín $\mathbb{A}^2 = \{x_1 = 0\} \setminus r$, debemos estudiar la restricción de $g$ al plano proyectivo $\pi \equiv \{x_1 = 0\}$ y encontrar una recta $r$ que, al ser tomada como recta del infinito, convierta la aplicación inducida en una homotecia. (Podemos interpretarlo tambien como que, la aplicación proyectiva debe ser una homología, es decir, que en alguna referencia debe tener la forma diag(1,k,k))
Consideramos el subespacio de puntos de la forma $(x_0 : 0 : x_2 : x_3)$. Aplicando las ecuaciones de la homografía original, obtenemos la matriz restringida $M$ que actúa sobre las coordenadas $(x_0, x_2, x_3)$:
\[
M = \begin{pmatrix} 
-1 & 0 & 0 \\ 
0 & -1 & 0 \\ 
0 & 3 & 2 
\end{pmatrix}\]
Una homología proyectiva es una proyectividad plana que posee una recta de puntos fijos (eje) y un punto fijo aislado (centro). Analizamos los autovalores de $M$:
\begin{itemize}
    \item \textbf{Eje de puntos fijos ($\lambda = -1$):} Resolvemos el sistema $(M + I)\vec{x} = \vec{0}$:
    \[
    \begin{pmatrix} 0 & 0 & 0 \\ 0 & 0 & 0 \\ 0 & 3 & 3 \end{pmatrix} \begin{pmatrix} x_0 \\ x_2 \\ x_3 \end{pmatrix} = \begin{pmatrix} 0 \\ 0 \\ 0 \end{pmatrix} \implies x_2 + x_3 = 0
    \]
    Todos los puntos de la recta $e \equiv \{x_2 + x_3 = 0\}$ dentro del plano $x_1=0$ son puntos fijos. Este es el \textbf{eje de la homología}.
    
    \item \textbf{Centro de la homología ($\lambda = 2$):} Resolvemos $(M - 2I)\vec{x} = \vec{0}$:
    \[
    \begin{pmatrix} -3 & 0 & 0 \\ 0 & -3 & 0 \\ 0 & 3 & 0 \end{pmatrix} \begin{pmatrix} x_0 \\ x_2 \\ x_3 \end{pmatrix} = \vec{0} \implies x_0 = 0, x_2 = 0
    \]
    Obtenemos el punto fijo aislado $C(0:0:1)$, que corresponde a $P_4(0:0:0:1)$ en $\mathbb{P}^3$.
\end{itemize}


Como proponemos, definimos una referencia proyectiva en $\pi$ formada por:
\begin{itemize}
    \item $u_1 = (0, 0, 1)$ (Centro $C$).
    \item $u_2 = (1, 0, 0)$ (Punto del eje $e$).
    \item $u_3 = (0, 1, -1)$ (Punto del eje $e$).
\end{itemize}
En esta referencia, la matriz de la aplicación es diagonal: $D = \text{diag}(2, -1, -1)$, que es proyectivamente equivalente a $\text{diag}(-2, 1, 1)$. 

Si elegimos como **recta del infinito** $r$ al propio eje de la homología ($r = e$), cualquier punto de dicha recta se transforma mediante la identidad (puntos fijos). En la carta afín resultante $\{x_1 = 0\} \setminus r$:
\begin{enumerate}
    \item La aplicación es una afinidad porque preserva la recta del infinito $r$.
    \item Al ser $r$ una recta de puntos fijos, la parte lineal es la identidad o un múltiplo de ella.
\end{enumerate}
e) La existencia de un punto fijo $C$ fuera del infinito garantiza que la aplicación sea una \textbf{homotecia} de centro $C$ (en nuestras coordenadas afines será el (0,0)) y razón $k = \frac{\lambda_\infty}{\lambda_C} = \frac{1}{-2}$
}

%ejercicioN(6)
\ejercicio{La polaridad asociada a una cónica no singular tiene las siguientes propiedades:
\begin{itemize}
    \item[(i)] La recta polar del punto $A_{1}=(1:0:-1)$ es $r_{1}:x_{0}+x_{1}-2x_{2}=0$.
    \item[(ii)] El polo de la recta $r_{2}:x_{2}=0$ es $A_{2}=(1:1:0)$ (esto es equivalente a que la polar de $A_{2}$ sea $r_2$).
    \item[(iii)] La recta $r_{3}:x_{0}=0$ es la polar de un punto de la cónica.
\end{itemize}
Se pide:\\
a) La recta polar del punto $A_{4}=(1:-1:0)$.\\
b) El polo de la recta $s:x_{0}+x_{1}+x_{2}=0$.\\
c) ¿Cuántas tangentes a la cónica pasan por el punto $A_{2}$?\\
d) Clasificar la cónica afín inducida en la carta afín $\{x_{2}\neq 0\}$.\\
e) La ecuación de la cónica.}{

Sea $\mathcal{C}$ una cónica no singular en $\mathbb{P}_{\mathbb{R}}^{2}$ asociada a una matriz simétrica $C$:
\[
C = \begin{pmatrix} 
c_{00} & c_{01} & c_{02} \\ 
c_{01} & c_{11} & c_{12} \\ 
c_{02} & c_{12} & c_{22} 
\end{pmatrix}
\]



\begin{itemize}
    \item \textbf{Propiedad (ii):} El polo de la recta $x_2=0$ (coeficientes $[0,0,1]$) es el punto $A_2(1:1:0)$. Por la relación fundamental de polaridad $C \cdot P = \rho L$:
    \[
    \begin{pmatrix} c_{00} & c_{01} & c_{02} \\ c_{01} & c_{11} & c_{12} \\ c_{02} & c_{12} & c_{22} \end{pmatrix} \begin{pmatrix} 1 \\ 1 \\ 0 \end{pmatrix} = \begin{pmatrix} 0 \\ 0 \\ \rho_2 \end{pmatrix} \implies 
    \begin{cases} 
    c_{00} + c_{01} = 0 \\ 
    c_{01} + c_{11} = 0 
    \end{cases}
    \]
    Tomando $c_{00} = k$ como factor de escala, obtenemos $c_{01} = -k$ y $c_{11} = k$.

    \item \textbf{Propiedad (i):} La recta polar de $A_1(1:0:-1)$ es $x_0 + x_1 - 2x_2 = 0$ (coeficientes $[1,1,-2]$):
    \[
    \begin{pmatrix} k & -k & c_{02} \\ -k & k & c_{12} \\ c_{02} & c_{12} & c_{22} \end{pmatrix} \begin{pmatrix} 1 \\ 0 \\ -1 \end{pmatrix} = \rho_1 \begin{pmatrix} 1 \\ 1 \\ -2 \end{pmatrix} \implies
    \begin{cases} 
    k - c_{02} = \rho_1 \implies c_{02} = k - \rho_1 \\ 
    -k - c_{12} = \rho_1 \implies c_{12} = -k - \rho_1 \\
    c_{02} - c_{22} = -2\rho_1 \implies (k - \rho_1) - c_{22} = -2\rho_1 \implies c_{22} = k + \rho_1
    \end{cases}
    \]
\end{itemize}

La matriz de la cónica queda definida en función de $k$ y el parámetro de proporcionalidad $\rho_1$:
\[
C = \begin{pmatrix} 
k & -k & k-\rho_1 \\ 
-k & k & -k-\rho_1 \\ 
k-\rho_1 & -k-\rho_1 & k+\rho_1 
\end{pmatrix}
\]



La propiedad (iii) establece que la recta $r_3: x_0 = 0$ es la polar de un punto de la cónica, lo cual implica que $r_3$ es una \textbf{recta tangente}. 

Geométricamente, la intersección de una cónica con una recta tangente debe ser un único punto doble. Consideramos la restricción de la cónica a la recta proyectiva $\mathbb{P}^1$ definida por $x_0 = 0$. Los puntos de esta recta son de la forma $(0:x_1:x_2)$. La ecuación de la cuádrica inducida en la recta es:
\[
\begin{pmatrix} x_1 & x_2 \end{pmatrix} 
\begin{pmatrix} c_{11} & c_{12} \\ c_{12} & c_{22} \end{pmatrix} 
\begin{pmatrix} x_1 \\ x_2 \end{pmatrix} = 0
\]
Según la clasificación de cuádricas en la recta proyectiva, para que existan un único punto doble (tangencia), el determinante de la matriz de restricción debe ser nulo:
\[
\det \begin{pmatrix} k & -k-\rho_1 \\ -k-\rho_1 & k+\rho_1 \end{pmatrix} = 0
\]
\[
k(k+\rho_1) - (-k-\rho_1)^2 = 0 \implies k^2 + k\rho_1 - (k^2 + 2k\rho_1 + \rho_1^2) = 0
\]
\[
-k\rho_1 - \rho_1^2 = 0 \implies -\rho_1(k + \rho_1) = 0
\]
Dado que la cónica es no singular ($\det C \neq 0$), descartamos $\rho_1 = 0$. Por tanto, obtenemos la condición $\rho_1 = -k$.



Sustituyendo $\rho_1 = -k$ en la matriz $C$ y fijando $k=1$ para normalizar:
\[
C = \begin{pmatrix} 
1 & -1 & 2 \\ 
-1 & 1 & 0 \\ 
2 & 0 & 0 
\end{pmatrix}
\]
a)
La recta polar $r$ de un punto $P$ se obtiene mediante la relación $L = C \cdot P$, donde $L$ es el vector de coeficientes de la recta $[u_0, u_1, u_2]$.
\[
L = \begin{pmatrix} 1 & -1 & 2 \\ -1 & 1 & 0 \\ 2 & 0 & 0 \end{pmatrix} \begin{pmatrix} 1 \\ -1 \\ 0 \end{pmatrix} = 
\begin{pmatrix} 1(1) + (-1)(-1) + 2(0) \\ -1(1) + 1(-1) + 0(0) \\ 2(1) + 0(-1) + 0(0) \end{pmatrix} = 
\begin{pmatrix} 2 \\ -2 \\ 2 \end{pmatrix}
\]
Dividiendo por el factor común $2$, obtenemos los coeficientes $[1, -1, 1]$. \\
\textbf{Ecuación de la recta polar:} $x_0 - x_1 + x_2 = 0$.

b) El polo de la recta $s : x_0 + x_1 + x_2 = 0$

Para hallar el polo $P(x_0:x_1:x_2)$, resolvemos el sistema $C \cdot P = \lambda \vec{s}$, donde $\vec{s} = [1, 1, 1]^T$:
\[
\begin{pmatrix} 1 & -1 & 2 \\ -1 & 1 & 0 \\ 2 & 0 & 0 \end{pmatrix} \begin{pmatrix} x_0 \\ x_1 \\ x_2 \end{pmatrix} = \begin{pmatrix} 1 \\ 1 \\ 1 \end{pmatrix} \implies
\begin{cases} 
x_0 - x_1 + 2x_2 = 1 \\ 
-x_0 + x_1 = 1 \\ 
2x_0 = 1 
\end{cases}
\]
1. De la tercera ecuación: $x_0 = 1/2$. \\
2. Sustituyendo en la segunda: $-1/2 + x_1 = 1 \implies x_1 = 3/2$. \\
3. Sustituyendo en la primera: $1/2 - 3/2 + 2x_2 = 1 \implies -1 + 2x_2 = 1 \implies x_2 = 1$. \\
El punto es $(1/2 : 3/2 : 1)$, que normalizado es $(1 : 3 : 2)$. \\
\textbf{Polo de la recta $s$:} $P = (1 : 3 : 2)$.

c) Cantidad de tangentes que pasan por $A_2 = (1 : 1 : 0)$

Evaluamos si el punto $A_2$ pertenece a la cónica usando su forma cuadrática $X^T C X = 0$:
\[
(1, 1, 0) \begin{pmatrix} 1 & -1 & 2 \\ -1 & 1 & 0 \\ 2 & 0 & 0 \end{pmatrix} \begin{pmatrix} 1 \\ 1 \\ 0 \end{pmatrix} = (0, 0, 2) \begin{pmatrix} 1 \\ 1 \\ 0 \end{pmatrix} = 0
\]
Como $A_2$ satisface la ecuación de la cónica, el punto es un punto de la curva. Por definición, en una cónica no singular, por cada punto de la curva pasa **exactamente una recta tangente** (su recta polar). \\
Pasa 1 tangente.

d) Clasificación de la cónica en la carta afín $\{x_2 \neq 0\}$

Intersecamos la cónica con la recta del infinito $x_2 = 0$:
\[
x_0^2 + x_1^2 - 2x_0x_1 + 4x_0(0) = 0 \implies (x_0 - x_1)^2 = 0
\]
La intersección produce un único punto real doble $(1:1:0)$. \\
 La cónica es una \textbf{parábola}.

e) Ecuación de la cónica

Expandiendo la expresión $X^T C X = 0$:
\[
\begin{pmatrix} x_0 & x_1 & x_2 \end{pmatrix} \begin{pmatrix} 1 & -1 & 2 \\ -1 & 1 & 0 \\ 2 & 0 & 0 \end{pmatrix} \begin{pmatrix} x_0 \\ x_1 \\ x_2 \end{pmatrix} = 0
\]
 $x_0^2 + x_1^2 - 2x_0x_1 + 4x_0x_2 = 0$.


}

%ejercicioN(7)
\ejercicio{Se tiene una familia uniparamétrica de cónicas no singulares $\{\varphi_{\alpha}\}$ cuyas polaridades asociadas tienen todas las siguientes propiedades:
\begin{itemize}
    \item[(i)] La recta polar del punto $A_{1}=(1:0:-1)$ pasa por $A_{1}$.
    \item[(ii)] El polo de la recta $r_{2}:x_{2}=0$ es $A_{2}=(0:1:0)$ (esto es equivalente a que la polar de $A_{2}$ sea $r_2$).
    \item[(iii)] La recta $r_{3}:-x_{0}+x_{1}=0$ es la polar del punto $A_{3}=(-1:1:1)$.
\end{itemize}
Se pide:\\
a) La recta polar de $A_{4}=(1:1:0)$.\\
b) Probar que la recta polar de $A_{1}$ es $x_{0}-x_{1}+x_{2}=0$.\\
c) Las ecuaciones de las cónicas de la familia (dependiendo de un parámetro $a$).\\
d) Clasificar las cónicas afines inducidas en $\{x_{0}+x_{1}\neq 0\}$.\\
e) Todas las rectas proyectivas $s$ que pasan por $A_{1}$ y en cuyo complementario alguna cónica de la familia induzca una parábola.}{
    Sea $C_{\alpha}$ la matriz simétrica asociada a la familia de cónicas $\{\varphi_{\alpha}\}$ en $\mathbb{P}_{\mathbb{R}}^{2}$. A partir de las propiedades de polaridad dadas:

\begin{enumerate}
    \item \textbf{Propiedad (ii):} El polo de la recta $x_2=0$ es $A_2(0:1:0)$. Esto implica que la segunda columna de la matriz es proporcional al vector de la recta $[0,0,1]^T$, lo que nos da $c_{01}=0, c_{11}=0$ y $c_{12} = \lambda$. Tomando $\lambda=1$:
    \[ C = \begin{pmatrix} c_{00} & 0 & c_{02} \\ 0 & 0 & 1 \\ c_{02} & 1 & c_{22} \end{pmatrix} \]
    
    \item \textbf{Propiedad (iii):} La polar de $A_3(-1:1:1)$ es $-x_0 + x_1 = 0$. Aplicando $C \cdot A_3 = \mu [-1, 1, 0]^T$:
    \[ \begin{pmatrix} c_{00} & 0 & c_{02} \\ 0 & 0 & 1 \\ c_{02} & 1 & c_{22} \end{pmatrix} \begin{pmatrix} -1 \\ 1 \\ 1 \end{pmatrix} = \begin{pmatrix} -c_{00} + c_{02} \\ 1 \\ -c_{02} + 1 + c_{22} \end{pmatrix} = \begin{pmatrix} -\mu \\ \mu \\ 0 \end{pmatrix} \]
    De la segunda fila, $\mu = 1$. Resolviendo: $c_{02} = c_{00} - 1$ y $c_{22} = c_{00} - 2$.
\end{enumerate}

Definiendo $\alpha = c_{00}$, la matriz de la familia es:
\[ C_{\alpha} = \begin{pmatrix} \alpha & 0 & \alpha-1 \\ 0 & 0 & 1 \\ \alpha-1 & 1 & \alpha-2 \end{pmatrix} \]

---



a) La recta polar de $A_4 = (1:1:0)$
$L_{A4} = C_{\alpha} \cdot (1, 1, 0)^T = (\alpha, 0, \alpha)^T$. Normalizando por $\alpha \neq 0$:
\textbf{Ecuación:} $x_0 + x_2 = 0$.

b) Probar que la recta polar de $A_1(1:0:-1)$ es $x_0 - x_1 + x_2 = 0$
$L_{A1} = C_{\alpha} \cdot (1, 0, -1)^T = (\alpha - (\alpha - 1), -1, (\alpha - 1) - (\alpha - 2))^T = (1, -1, 1)^T$.
La ecuación es $x_0 - x_1 + x_2 = 0$. Puesto que $1 - 0 + (-1) = 0$, se cumple que el polo $A_1$ pertenece a su propia polar (Propiedad i).

c) Ecuaciones de las cónicas de la familia
Ecuación: $\alpha x_0^2 + 2(\alpha-1)x_0x_2 + 2x_1x_2 + (\alpha-2)x_2^2 = 0$.

d) Clasificación en la carta afín $\{x_0 + x_1 \neq 0\}$

Primero, verificamos la \textbf{no singularidad} de la familia. La cónica es no singular si su determinante es no nulo:
\[ |C_{\alpha}| = \det \begin{pmatrix} \alpha & 0 & \alpha-1 \\ 0 & 0 & 1 \\ \alpha-1 & 1 & \alpha-2 \end{pmatrix} = -1 \cdot \det \begin{pmatrix} \alpha & 0 \\ \alpha-1 & 1 \end{pmatrix} = -\alpha \]
Por tanto, la cónica es \textbf{no singular si y solo si $\alpha \neq 0$}.

Para clasificar, realizamos el cambio de referencia para que la recta del infinito sea $X_0 = x_0 + x_1$:
Definimos $X_0 = x_0 + x_1, X_1 = x_1, X_2 = x_2$. La matriz de paso es $M = \begin{pmatrix} 1 & -1 & 0 \\ 0 & 1 & 0 \\ 0 & 0 & 1 \end{pmatrix}$.
La matriz en la nueva referencia es $C' = M^T C_{\alpha} M$:
\[ C' = \begin{pmatrix} \alpha & -\alpha & \alpha-1 \\ -\alpha & \alpha & 2-\alpha \\ \alpha-1 & 2-\alpha & \alpha-2 \end{pmatrix} \]
El tipo afín depende de la submatriz de direcciones infinitas $A_{00}$ (filas y columnas 1 y 2):
\[ |A_{00}| = \det \begin{pmatrix} \alpha & 2-\alpha \\ 2-\alpha & \alpha-2 \end{pmatrix} = 2\alpha - 4 \]
Clasificación ($\alpha \neq 0$):
\begin{itemize}
    \item \textbf{Elipse:} $2\alpha - 4 > 0 \implies \alpha > 2$.
    \item \textbf{Parábola:} $2\alpha - 4 = 0 \implies \alpha = 2$.
    \item \textbf{Hipérbola:} $2\alpha - 4 < 0 \implies \alpha < 2$ (con $\alpha \neq 0$).
\end{itemize}

e) Rectas $s$ por $A_1$ que inducen una parábola
Una cónica induce una parábola en el complementario de una recta $s$ si y solo si $s$ es tangente a la cónica. Como $A_1$ es un punto de todas las cónicas de la familia, la recta $s$ debe ser la tangente en $A_1$, que es su polar.
Del apartado (b), la polar de $A_1$ es única e independiente de $\alpha$:
\textbf{Respuesta:} $s: x_0 - x_1 + x_2 = 0$.
}

%ejercicioN(8)
\ejercicio{Decide razonadamente para qué cónicas existen referencias proyectivas tales que su ecuación se escribe en coordenadas como $x_{0}^{2}-x_{1}x_{2}=0$ (distingue si es necesario el caso real y el complejo). Describe geométricamente la posición de los puntos que forman parte de la referencia en función del lugar de ceros de la cónica.}{
    Dada la ecuación de la cónica $C: x_0^2 - x_1x_2 = 0$, podemos escribir su matriz asociada $A$ en la referencia proyectiva dada:
\[
A = \begin{pmatrix} 
1 & 0 & 0 \\ 
0 & 0 & -1/2 \\ 
0 & -1/2 & 0 
\end{pmatrix}
\]

Para determinar el tipo de cónica, calculamos el determinante de la matriz:
\[
\det(A) = 1 \cdot \left( 0 - \left(-\frac{1}{2}\right)\left(-\frac{1}{2}\right) \right) = -\frac{1}{4}
\]

Como $\det(A) \neq 0$, la cónica es \textbf{no degenerada} (rango 3). 
Caso Complejo ($\mathbb{P}^2(\mathbb{C})$)
En el plano proyectivo complejo, todas las cónicas no degeneradas son proyectivamente equivalentes. Por tanto, existe tal referencia para \textbf{cualquier cónica propia}.

Caso Real ($\mathbb{P}^2(\mathbb{R})$)
En el caso real, debemos considerar la signatura de la forma cuadrática. Los autovalores de $A$ son $\{1, 1/2, -1/2\}$, lo que da una signatura $(2, 1)$. Esto corresponde a una \textbf{cónica real con puntos} (no vacía). En términos afines, esto representa elipses, parábolas o hipérbolas reales.

Descripción geométrica de la referencia

Sea la referencia $\mathcal{R} = \{P_0, P_1, P_2; U\}$. Analizamos la posición de estos puntos respecto a la cónica $C$ mediante sus coordenadas:

\begin{itemize}
    \item \textbf{Puntos en la cónica:} 
    Sustituyendo las coordenadas en la ecuación $x_0^2 - x_1x_2 = 0$:
    \begin{itemize}
        \item $P_1 = [0 : 1 : 0] \implies 0^2 - (1)(0) = 0 \quad \therefore P_1 \in C$
        \item $P_2 = [0 : 0 : 1] \implies 0^2 - (0)(1) = 0 \quad \therefore P_2 \in C$
        \item $U = [1 : 1 : 1] \implies 1^2 - (1)(1) = 0 \quad \therefore U \in C$
    \end{itemize}
    
    \item \textbf{Punto fuera de la cónica:}
    \begin{itemize}
        \item $P_0 = [1 : 0 : 0] \implies 1^2 - (0)(0) = 1 \neq 0 \quad \therefore P_0 \notin C$
    \end{itemize}
\end{itemize}


La recta polar de un punto $P$ respecto a la cónica se define como $P^T A X = 0$.

\begin{enumerate}
    \item \textbf{Tangente en $P_1$:} La polar de $P_1$ es $x_2 = 0$. Esta es la recta que une $P_0$ y $P_1$. Por lo tanto, la recta $\overline{P_0P_1}$ es la recta tangente a la cónica en el punto $P_1$.
    \item \textbf{Tangente en $P_2$:} La polar de $P_2$ es $x_1 = 0$. Esta es la recta que une $P_0$ y $P_2$. Por lo tanto, la recta $\overline{P_0P_2}$ es la recta tangente a la cónica en el punto $P_2$.
\end{enumerate}


Para que una cónica pueda expresarse de esta forma, debe ser \textbf{no degenerada} (y no vacía en el caso real). Los elementos de la referencia se sitúan de la siguiente manera:
\begin{itemize}
    \item $P_1$ y $P_2$ son dos puntos de la propia cónica.
    \item $P_0$ es el punto donde se cruzan las rectas tangentes a la cónica que pasan por $P_1$ y $P_2$.
    \item $U$ (punto unidad) es un tercer punto de la cónica que fija la escala de la referencia.
\end{itemize}
}

%ejercicioN(9)
\ejercicio{Dada una cuádrica proyectiva, sea $r$ una recta que no sea tangente ni pase por un punto singular.\\
a) Muestra que $g:r\rightarrow r$ definida por $g(A)=Polar_{\varphi}(A)\cap r$ está bien definida y es una homografía.\\
b) Halla los puntos fijos de $g$ y muestra que $g\circ g=id$ (en otras palabras, $g$ es involutiva).\\
c) Suponiendo que $r\cap\mathcal{C}_{\varphi}=\{B,B^{\prime}\}$, prueba que para todo $A\in r \setminus \{B, B^{\prime}\}$ se cumple $[B,B^{\prime},A,g(A)]=-1$.}{
    Sea $\varphi$ una cuádrica proyectiva y $r$ una recta que no es tangente a $\varphi$ ni pasa por un punto singular.

a) 

Sea $A \in r$. Bajo la polaridad inducida por la cuádrica $\varphi$, el punto $A$ se transforma en su hiperplano polar $\pi_A$ en el espacio dual. Para que la aplicación $g(A) = \pi_A \cap r$ esté bien definida, la intersección debe ser un punto único de la recta, lo que exige que $r \not\subset \pi_A$ (de lo contrario, la intersección sería toda la recta $r$).

Razonamos por reducción al absurdo: supongamos que $r \subseteq \pi_A$. Como $A \in r$, esto implicaría que $A \in \pi_A$, lo que por definición significa que $A$ es un punto de la cuádrica ($A \in \mathcal{C}_\varphi$). Al ser $A$ un punto de la cuádrica, su hiperplano polar $\pi_A$ es el hiperplano tangente en $A$. Si la recta $r$ estuviera contenida en el hiperplano tangente y pasara por el punto de contacto $A$, entonces $r$ sería una recta tangente a la cuádrica. Esto supone una \textbf{contradicción} con la hipótesis inicial.

Por tanto, podemos afirmar que para ningún punto de la recta los hiperplanos polares contienen a $r$. Así, la recta no está contenida en la base del haz y podemos definir $g$ como la composición de dos aplicaciones proyectivas: la polaridad $f: r \to \text{Polar}_\varphi(r)$ y la intersección con la recta $r$. Al ser composición de homografías, $g$ es una homografía.

b) 

Al ser $g$ una homografía, un punto es fijo si $g(x) = x$. Esto ocurre si y solo si $x$ pertenece a su propio hiperplano polar ($x \in \pi_x$), lo cual es la condición necesaria y suficiente para que $x$ sea un punto de la cuádrica ($x \in \mathcal{C}_\varphi$). Por tanto, los puntos fijos de $g$ son exactamente los puntos de la intersección de la recta con la cuádrica:
\[ \text{Fix}(g) = r \cap \mathcal{C}_\varphi \]

Demostrar que $g$ es una involución es inmediato por la reprocidad de los polos y polares: si el punto $A'$ está en el polar de $A$, entonces el punto $A$ debe estar en el polar de $A'$. Analíticamente, si $\varphi(A, A') = 0$, entonces $\varphi(A', A) = 0$, lo que implica que $g(g(A)) = A$.

c) 

Sea $r \cap \mathcal{C}_\varphi = \{B, B'\}$. Para demostrar que la razón doble es $-1$, elegimos una referencia proyectiva sobre la recta $r$ utilizando los puntos fijos de la involución:
\begin{itemize}
    \item $B = [1:0]$ (origen).
    \item $B' = [0:1]$ (punto del infinito).
\end{itemize}

Cualquier punto $A \in r$ puede expresarse como $A = B + x B'$, donde $x$ es su coordenada. Sea $g(A) = B + x' B'$ su imagen. La condición de que $A$ y $g(A)$ sean conjugados respecto a la cuádrica $\varphi$ (con matriz $M$) es $A^T M g(A) = 0$. Desarrollando:
\[ (B + xB')^T M (B + x'B') = 0 \]
\[ B^T MB + x'(B^T MB') + x(B'^T MB) + xx'(B'^T MB') = 0 \]

Como $B$ y $B'$ pertenecen a la cuádrica, $B^T MB = 0$ y $B'^T MB' = 0$. Por simetría, $B^T MB' = B'^T MB = k$. La ecuación se simplifica a:
\[ kx' + kx = 0 \implies x' = -x \]

La homografía en esta referencia es la aplicación $x \mapsto -x$. Calculamos la razón doble de los cuatro puntos utilizando sus coordenadas $\{0, \infty, x, -x\}$:
\[ [B, B', A, g(A)] = \frac{x - 0}{x - \infty} : \frac{-x - 0}{-x - \infty} \]
En el límite cuando una coordenada tiende a infinito, la expresión se reduce a:
\[ [0, \infty, x, -x] = \frac{x}{-x} = -1 \]
Esto demuestra que los puntos $A$ y $g(A)$ están en división armónica respecto a los puntos de intersección de la recta con la cuádrica.
}

%ejercicioN(10)
\ejercicio{Prueba que la dimensión del espacio proyectivo de las cuádricas de $\mathbb{P}_{\mathbb{R}}^{3}$ es 9, y deduce que cualesquiera tres rectas de $\mathbb{P}_{\mathbb{R}}^{3}$ están contenidas en (el lugar de ceros de) una misma cuádrica.\\
$\bullet$ Estudia la posición relativa de tres rectas contenidas en una cuádrica degenerada en $\mathbb{P}_{\mathbb{R}}^{3}$. Elabora una lista de posibilidades en función del tipo de cuádrica.\\
$\bullet$ Prueba o refuta la siguiente afirmación: dadas tres rectas disjuntas en $\mathbb{P}_{\mathbb{R}}^{3}$, no hay ninguna cuádrica degenerada que las contenga.}{}

%ejercicioN(11)
\ejercicio{Sobre superficies cuádricas en dimensión 3:\\
a) Demuestra que la cuádrica proyectiva real de signatura $(3,1)$ no contiene rectas en su lugar de ceros (y de ahí su denominación cuádrica no reglada).\\
b) Dado un punto $p$ cualquiera de la cuádrica proyectiva real de signatura $(2,2)$, obtén una recta que pasa por $p$ y está contenida en la cuádrica. ¿Hay más de una? Deduce que (el lugar de ceros de) la cuádrica es unión de rectas, por eso se denomina cuádrica reglada.\\
c) Comprueba que la superficie cuádrica compleja no degenerada es también una cuádrica reglada: su lugar de ceros es una unión de rectas.\\
d) Utiliza los dos primeros apartados para enumerar cuáles son las posibles cónicas obtenidas por restricción de la cuádrica reglada y la no reglada a un plano (y así justificar las posibles cónicas de infinito estudiadas en los casos de clasificación de superficies cuádricas afines vistos en clase).}{}

%ejercicioN(12)
\ejercicio{Sea una cuádrica no degenerada, $p\in\mathcal{C}_{\varphi}$ y $H$ un hiperplano que no pasa por $p$.\\
a) Decide para qué puntos $q$ la asignación $q\mapsto h(q)=V(p,q)\cap(\mathcal{C}_{\varphi}\setminus\{p\})$ está bien definida.\\
b) Determina $X$ de forma que $h:H\setminus X\rightarrow\mathcal{C}_{\varphi}\setminus\{p\}$ es una biyección (que se llama proyección estereográfica).\\
c) Supongamos que es una cuádrica en $P_{\mathbb{K}}^{n}$, siendo $\mathbb{K}=\mathbb{Z}_{m}$ con $m$ primo el cuerpo con $m$ elementos. Halla el número de puntos que contiene $H\setminus X$ en función de $m$ y $n$ y utiliza la proyección estereográfica para dar una cota inferior al número de puntos contenidos en una cuádrica en $P_{\mathbb{K}}^{n}$ supuesta que es no vacía.\\
d) Los primeros apartados proporcionan un método para parametrizar cónicas: se considera una parametrización $t\mapsto\alpha(t)$ de una recta afín que no pase por $p$ (que hará las veces de $H\setminus X$) y la parametrización buscada es $t\mapsto h(\alpha(t))$. Parametriza $x_{0}^{2}+x_{0}x_{1}-2x_{1}^{2}+9x_{2}^{2}=0$ desde el punto $p=(1:-2:1)$.}{}